% Expanded (post-lowering) form of arrus.tx_rx_sequence.
%
% Fig 6 shows the AUTHORED form: one region body describing one slot's
% schema, compact but hiding the fact that at runtime this is a
% sequence of side-effecting firings threaded through device state.
%
% This figure shows the EXPANDED form after Us4RBackend.compile:
%   - the region body's pure ops (aperture_for_slot, compute_delays)
%     are hoisted into compile-time-materialized LUTs;
%   - the region body's side-effecting op (tx_rx.slot) is expanded
%     into 192 concrete slot descriptors resident in device DDR;
%   - a runtime `us4r.rebind_wiring(@indices)` step updates the linked
%     list to wire only the selected slots into the active chain;
%   - `us4r.fire_sequence` threads a `!us4r.device_state` value through
%     ONLY the wired slots, producing the (K, ...) stack of RF frames.
%
% Selected example: @indices = [3, 5, 8, 13].
% Slot 3, 5, 8, 13 are HIGHLIGHTED as active; the others are grayed
% out to show they exist in device memory but are skipped this frame.
%
% Build: pdflatex ndarray_fig7_txrx_expanded.tex

\documentclass[tikz,border=6pt]{standalone}
% Shared preamble for NdArray graph figures.
% Include with: % Shared preamble for NdArray graph figures.
% Include with: % Shared preamble for NdArray graph figures.
% Include with: \input{ndarray_common}
%
% Provides a consistent color palette and node styles for the three figures
% in docs/td/ndarray_design.md.

\usepackage{tikz}
\usetikzlibrary{positioning, arrows.meta, fit, calc, backgrounds, shapes.geometric}

% --- Device / role color palette ------------------------------------------------
\definecolor{arrusUs4R}   {HTML}{2CA9BC}   % ultrasound device (teal)
\definecolor{arrusGPU}    {HTML}{F28C28}   % CUDA GPU (orange)
\definecolor{arrusCPU}    {HTML}{4CAF50}   % CPU (green)
\definecolor{arrusHost}   {HTML}{9E9E9E}   % host / control (gray)
\definecolor{arrusBuffer} {HTML}{616161}   % stream buffer (dark gray)
\definecolor{arrusVariable}{HTML}{F4C542}  % Variable (gold)
\definecolor{arrusTrigger}{HTML}{D64545}   % Trigger source (red)
\definecolor{arrusOutput} {HTML}{5C6BC0}   % Graph output / sink (indigo)

% --- Node styles ---------------------------------------------------------------
\tikzset{
  op/.style       = {rectangle, rounded corners=2pt, draw=black!60, thick,
                     minimum width=26mm, minimum height=8mm, align=center,
                     font=\small},
  op/us4r/.style  = {op, fill=arrusUs4R!25,   draw=arrusUs4R!80!black},
  op/gpu/.style   = {op, fill=arrusGPU!25,    draw=arrusGPU!80!black},
  op/cpu/.style   = {op, fill=arrusCPU!25,    draw=arrusCPU!70!black},
  op/host/.style  = {op, fill=arrusHost!20,   draw=arrusHost!70!black},
  op/output/.style= {op, fill=arrusOutput!25, draw=arrusOutput!80!black},
  op/variable/.style = {ellipse, draw=arrusVariable!70!black, thick,
                        fill=arrusVariable!30, minimum width=22mm,
                        minimum height=7mm, align=center, font=\small\itshape},
  op/trigger/.style  = {regular polygon, regular polygon sides=3,
                        draw=arrusTrigger!80!black, thick, fill=arrusTrigger!25,
                        minimum size=11mm, align=center, font=\scriptsize,
                        inner sep=0pt},
  op/buffer/.style   = {cylinder, shape border rotate=90, aspect=0.3,
                        draw=arrusBuffer!80!black, thick, fill=arrusBuffer!15,
                        minimum width=18mm, minimum height=10mm, align=center,
                        font=\scriptsize},
  edge/.style        = {-{Latex[length=2mm]}, thick, draw=black!70},
  % Intra-region data flow: shared buffer + CUDA-stream / kernel-order
  % ordering. No runtime queue.
  edge/intra/.style  = {-{Latex[length=2mm]}, thick, draw=black!70},
  % Cross-region stream: real queue infrastructure (backpressure/drop policy,
  % scheduler wake-up). Used for cross-device transfers and rate mismatches.
  edge/stream/.style = {-{Latex[length=2mm]}, very thick, draw=black!70,
                        double distance=1pt},
  edge/trigger/.style= {-{Latex[length=2mm]}, thick, draw=arrusTrigger!80!black,
                        dashed},
  edge/variable/.style = {-{Latex[length=2mm]}, thick,
                          draw=arrusVariable!80!black, densely dotted},
  region/.style      = {rounded corners, inner sep=4mm, dashed,
                        line width=0.8pt},
  region/us4r/.style = {region, draw=arrusUs4R!70!black, fill=arrusUs4R!5},
  region/gpu/.style  = {region, draw=arrusGPU!70!black,  fill=arrusGPU!5},
  legendbox/.style   = {rectangle, draw=black!30, rounded corners=2pt,
                        inner sep=4pt, fill=white},
}

%
% Provides a consistent color palette and node styles for the three figures
% in docs/td/ndarray_design.md.

\usepackage{tikz}
\usetikzlibrary{positioning, arrows.meta, fit, calc, backgrounds, shapes.geometric}

% --- Device / role color palette ------------------------------------------------
\definecolor{arrusUs4R}   {HTML}{2CA9BC}   % ultrasound device (teal)
\definecolor{arrusGPU}    {HTML}{F28C28}   % CUDA GPU (orange)
\definecolor{arrusCPU}    {HTML}{4CAF50}   % CPU (green)
\definecolor{arrusHost}   {HTML}{9E9E9E}   % host / control (gray)
\definecolor{arrusBuffer} {HTML}{616161}   % stream buffer (dark gray)
\definecolor{arrusVariable}{HTML}{F4C542}  % Variable (gold)
\definecolor{arrusTrigger}{HTML}{D64545}   % Trigger source (red)
\definecolor{arrusOutput} {HTML}{5C6BC0}   % Graph output / sink (indigo)

% --- Node styles ---------------------------------------------------------------
\tikzset{
  op/.style       = {rectangle, rounded corners=2pt, draw=black!60, thick,
                     minimum width=26mm, minimum height=8mm, align=center,
                     font=\small},
  op/us4r/.style  = {op, fill=arrusUs4R!25,   draw=arrusUs4R!80!black},
  op/gpu/.style   = {op, fill=arrusGPU!25,    draw=arrusGPU!80!black},
  op/cpu/.style   = {op, fill=arrusCPU!25,    draw=arrusCPU!70!black},
  op/host/.style  = {op, fill=arrusHost!20,   draw=arrusHost!70!black},
  op/output/.style= {op, fill=arrusOutput!25, draw=arrusOutput!80!black},
  op/variable/.style = {ellipse, draw=arrusVariable!70!black, thick,
                        fill=arrusVariable!30, minimum width=22mm,
                        minimum height=7mm, align=center, font=\small\itshape},
  op/trigger/.style  = {regular polygon, regular polygon sides=3,
                        draw=arrusTrigger!80!black, thick, fill=arrusTrigger!25,
                        minimum size=11mm, align=center, font=\scriptsize,
                        inner sep=0pt},
  op/buffer/.style   = {cylinder, shape border rotate=90, aspect=0.3,
                        draw=arrusBuffer!80!black, thick, fill=arrusBuffer!15,
                        minimum width=18mm, minimum height=10mm, align=center,
                        font=\scriptsize},
  edge/.style        = {-{Latex[length=2mm]}, thick, draw=black!70},
  % Intra-region data flow: shared buffer + CUDA-stream / kernel-order
  % ordering. No runtime queue.
  edge/intra/.style  = {-{Latex[length=2mm]}, thick, draw=black!70},
  % Cross-region stream: real queue infrastructure (backpressure/drop policy,
  % scheduler wake-up). Used for cross-device transfers and rate mismatches.
  edge/stream/.style = {-{Latex[length=2mm]}, very thick, draw=black!70,
                        double distance=1pt},
  edge/trigger/.style= {-{Latex[length=2mm]}, thick, draw=arrusTrigger!80!black,
                        dashed},
  edge/variable/.style = {-{Latex[length=2mm]}, thick,
                          draw=arrusVariable!80!black, densely dotted},
  region/.style      = {rounded corners, inner sep=4mm, dashed,
                        line width=0.8pt},
  region/us4r/.style = {region, draw=arrusUs4R!70!black, fill=arrusUs4R!5},
  region/gpu/.style  = {region, draw=arrusGPU!70!black,  fill=arrusGPU!5},
  legendbox/.style   = {rectangle, draw=black!30, rounded corners=2pt,
                        inner sep=4pt, fill=white},
}

%
% Provides a consistent color palette and node styles for the three figures
% in docs/td/ndarray_design.md.

\usepackage{tikz}
\usetikzlibrary{positioning, arrows.meta, fit, calc, backgrounds, shapes.geometric}

% --- Device / role color palette ------------------------------------------------
\definecolor{arrusUs4R}   {HTML}{2CA9BC}   % ultrasound device (teal)
\definecolor{arrusGPU}    {HTML}{F28C28}   % CUDA GPU (orange)
\definecolor{arrusCPU}    {HTML}{4CAF50}   % CPU (green)
\definecolor{arrusHost}   {HTML}{9E9E9E}   % host / control (gray)
\definecolor{arrusBuffer} {HTML}{616161}   % stream buffer (dark gray)
\definecolor{arrusVariable}{HTML}{F4C542}  % Variable (gold)
\definecolor{arrusTrigger}{HTML}{D64545}   % Trigger source (red)
\definecolor{arrusOutput} {HTML}{5C6BC0}   % Graph output / sink (indigo)

% --- Node styles ---------------------------------------------------------------
\tikzset{
  op/.style       = {rectangle, rounded corners=2pt, draw=black!60, thick,
                     minimum width=26mm, minimum height=8mm, align=center,
                     font=\small},
  op/us4r/.style  = {op, fill=arrusUs4R!25,   draw=arrusUs4R!80!black},
  op/gpu/.style   = {op, fill=arrusGPU!25,    draw=arrusGPU!80!black},
  op/cpu/.style   = {op, fill=arrusCPU!25,    draw=arrusCPU!70!black},
  op/host/.style  = {op, fill=arrusHost!20,   draw=arrusHost!70!black},
  op/output/.style= {op, fill=arrusOutput!25, draw=arrusOutput!80!black},
  op/variable/.style = {ellipse, draw=arrusVariable!70!black, thick,
                        fill=arrusVariable!30, minimum width=22mm,
                        minimum height=7mm, align=center, font=\small\itshape},
  op/trigger/.style  = {regular polygon, regular polygon sides=3,
                        draw=arrusTrigger!80!black, thick, fill=arrusTrigger!25,
                        minimum size=11mm, align=center, font=\scriptsize,
                        inner sep=0pt},
  op/buffer/.style   = {cylinder, shape border rotate=90, aspect=0.3,
                        draw=arrusBuffer!80!black, thick, fill=arrusBuffer!15,
                        minimum width=18mm, minimum height=10mm, align=center,
                        font=\scriptsize},
  edge/.style        = {-{Latex[length=2mm]}, thick, draw=black!70},
  % Intra-region data flow: shared buffer + CUDA-stream / kernel-order
  % ordering. No runtime queue.
  edge/intra/.style  = {-{Latex[length=2mm]}, thick, draw=black!70},
  % Cross-region stream: real queue infrastructure (backpressure/drop policy,
  % scheduler wake-up). Used for cross-device transfers and rate mismatches.
  edge/stream/.style = {-{Latex[length=2mm]}, very thick, draw=black!70,
                        double distance=1pt},
  edge/trigger/.style= {-{Latex[length=2mm]}, thick, draw=arrusTrigger!80!black,
                        dashed},
  edge/variable/.style = {-{Latex[length=2mm]}, thick,
                          draw=arrusVariable!80!black, densely dotted},
  region/.style      = {rounded corners, inner sep=4mm, dashed,
                        line width=0.8pt},
  region/us4r/.style = {region, draw=arrusUs4R!70!black, fill=arrusUs4R!5},
  region/gpu/.style  = {region, draw=arrusGPU!70!black,  fill=arrusGPU!5},
  legendbox/.style   = {rectangle, draw=black!30, rounded corners=2pt,
                        inner sep=4pt, fill=white},
}


\begin{document}
\begin{tikzpicture}[
    node distance=6mm and 8mm,
    every node/.append style={font=\small},
    slot/active/.style   = {rectangle, rounded corners=2pt, draw=arrusUs4R!80!black,
                            fill=arrusUs4R!30, thick, minimum width=16mm,
                            minimum height=14mm, align=center,
                            font=\scriptsize\bfseries},
    slot/inactive/.style = {rectangle, rounded corners=2pt, draw=black!30,
                            fill=black!5, minimum width=16mm, minimum height=14mm,
                            align=center, font=\scriptsize, text=black!45},
    op/const/.style      = {op, fill=arrusHost!20, draw=arrusHost!80!black,
                            font=\scriptsize\itshape},
    edge/data/.style     = {-{Latex[length=2mm]}, thick, draw=black!70},
    edge/state/.style    = {-{Latex[length=2.4mm]}, very thick,
                            draw=arrusTrigger!80!black},
    edge/skip/.style     = {-{Latex[length=1.5mm]}, thin, dashed, draw=black!35},
]

  % --- Top: @indices variable + LUT constants ------------------------------
  \node[op/variable, minimum width=52mm] (indices)
      {\textbf{Variable(@indices) = [3, 5, 8, 13]}\\%
       \scriptsize\itshape (K, i32) @host};

  \node[op/const, right=15mm of indices, minimum width=68mm] (luts)
      {\textbf{compile-time LUTs}\\%
       apertures\_LUT : const<192 $\times$ aperture> @Us4R:0\\%
       delays\_LUT   : const<192 $\times$ delays>    @Us4R:0\\%
       \scriptsize (hoisted from the region body's pure ops)};

  % --- Middle: rebind_wiring op --------------------------------------------
  \node[op/us4r, below=12mm of indices, minimum width=60mm] (rebind)
      {\textbf{us4r.rebind\_wiring}\\%
       \scriptsize (@indices) $\to$ \%wiring : !us4r.chain};

  % --- Slot strip (12 slots + two "..." breaks; 4 are active) --------------
  % Position the strip below with a horizontal layout.
  \node[slot/inactive] (slot0) at (-70mm, -80mm) {slot 0\\\scriptsize skipped};
  \node[slot/inactive, right=of slot0] (slot1) {slot 1\\\scriptsize skipped};
  \node[slot/inactive, right=of slot1] (slot2) {slot 2\\\scriptsize skipped};
  \node[slot/active,   right=of slot2] (slot3) {\textbf{slot 3}\\\scriptsize \textbf{fires}};
  \node[slot/inactive, right=of slot3] (slot4) {slot 4\\\scriptsize skipped};
  \node[slot/active,   right=of slot4] (slot5) {\textbf{slot 5}\\\scriptsize \textbf{fires}};
  \node[right=of slot5, font=\Large\bfseries] (dots1) {\ldots};
  \node[slot/active, right=of dots1] (slot8) {\textbf{slot 8}\\\scriptsize \textbf{fires}};
  \node[right=of slot8, font=\Large\bfseries] (dots2) {\ldots};
  \node[slot/active, right=of dots2] (slot13) {\textbf{slot 13}\\\scriptsize \textbf{fires}};
  \node[right=of slot13, font=\Large\bfseries] (dots3) {\ldots};
  \node[slot/inactive, right=of dots3] (slot191) {slot 191\\\scriptsize skipped};

  % Label for the strip
  \node[above=6mm of slot4, font=\small\itshape, text=arrusUs4R!80!black,
        align=center]
       (striplabel)
       {\textbf{us4r.fire\_sequence} \; --- 192 slot descriptors materialized in device DDR;\\%
        \%wiring selects the sub-chain to walk};

  % --- Initial device state (left of strip) --------------------------------
  \node[op/variable, left=8mm of slot0, minimum width=20mm,
        fill=arrusTrigger!25, draw=arrusTrigger!80!black]
       (s0) {\%s0\\\scriptsize\itshape !us4r.device\_state};

  % --- Final state + frames output (right of strip) ------------------------
  \node[op/output, right=10mm of slot191, minimum width=24mm] (frames)
      {\%frames\\\scriptsize (K, n\_ch, n\_samp)};
  \node[below=3mm of frames, op/variable, minimum width=24mm,
        fill=arrusTrigger!25, draw=arrusTrigger!80!black]
       (sfinal) {\%s1\\\scriptsize\itshape !us4r.device\_state};

  % --- Edges: variable / LUT reads ----------------------------------------
  \draw[edge/variable] (indices.south) -- (rebind.north)
      node[midway, right, font=\scriptsize\itshape,
           text=arrusVariable!80!black] {@indices};

  \draw[edge/data] (luts.south) |- ($(rebind.east)+(6mm,0)$) -- (rebind.east)
      node[near start, right, font=\scriptsize\itshape, text=black!55]
      {LUTs (baked-in per slot)};

  % rebind -> the strip (wiring)
  \draw[edge/data, rounded corners=6pt] (rebind.south) --
      ++(0,-6mm) node[midway, right, font=\scriptsize\itshape, text=black!55]
      {\%wiring} -| (striplabel.north);

  % --- State-thread edges (only through active slots, following wiring) ---
  \draw[edge/state] (s0.east)   -- (slot3.west)
      node[midway, above, font=\scriptsize\itshape,
           text=arrusTrigger!80!black] {\%s0};
  \draw[edge/state] (slot3.east) to[bend right=25] (slot5.west);
  \draw[edge/state] (slot5.east) to[bend right=20] (slot8.west);
  \draw[edge/state] (slot8.east) to[bend right=20] (slot13.west);
  \draw[edge/state] (slot13.east) to[bend right=15] (sfinal.west);

  % --- Skipped edges (dashed, showing linked-list bypass) -----------------
  \draw[edge/skip] (slot0.east) -- (slot1.west);
  \draw[edge/skip] (slot1.east) -- (slot2.west);
  \draw[edge/skip] (slot2.east) -- (slot3.west);
  \draw[edge/skip] (slot3.east) -- (slot4.west);
  \draw[edge/skip] (slot4.east) -- (slot5.west);

  % --- Frames output: each active slot contributes one row ---------------
  \draw[edge/data, rounded corners=6pt] (slot3.south)  |- ($(frames.west)+(-6mm,-6mm)$);
  \draw[edge/data, rounded corners=6pt] (slot5.south)  |- ($(frames.west)+(-6mm,-4mm)$);
  \draw[edge/data, rounded corners=6pt] (slot8.south)  |- ($(frames.west)+(-6mm,-2mm)$);
  \draw[edge/data, rounded corners=6pt] (slot13.south) |- (frames.west)
      node[near end, above, font=\scriptsize\itshape, text=black!55]
      {\%rf per active slot};

  % --- Legend --------------------------------------------------------------
  \node[legendbox, below=20mm of slot4, anchor=north] {%
    \begin{tikzpicture}[node distance=3mm]
      \node[slot/active, minimum width=26mm, minimum height=8mm] (l1)
        {active slot (in @indices)};
      \node[slot/inactive, minimum width=26mm, minimum height=8mm, below=of l1] (l2)
        {skipped slot};
      \node[op/const, minimum width=26mm, minimum height=6mm, below=of l2] (l3)
        {compile-time LUT};
      \node[below=3mm of l3.south, anchor=north, font=\scriptsize] (e1)
        {\tikz\draw[edge/state] (0,0) -- (10mm,0);\ !us4r.device\_state (state thread)};
      \node[below=1mm of e1.south, anchor=north, font=\scriptsize] (e2)
        {\tikz\draw[edge/data]  (0,0) -- (10mm,0);\ SSA data};
      \node[below=1mm of e2.south, anchor=north, font=\scriptsize] (e3)
        {\tikz\draw[edge/skip]  (0,0) -- (10mm,0);\ linked-list bypass (slot skipped)};
      \node[below=1mm of e3.south, anchor=north, font=\scriptsize] (e4)
        {\tikz\draw[edge/variable] (0,0) -- (10mm,0);\ variable read};
    \end{tikzpicture}%
  };

\end{tikzpicture}
\end{document}
